\documentclass[runningheads]{llncs}

\usepackage{graphicx}
\usepackage{booktabs}
\usepackage{amsmath}
\usepackage{hyperref}
\usepackage{caption}
\usepackage{subcaption}
\usepackage[table]{xcolor}

\definecolor{bestgreen}{RGB}{0,150,0}
\definecolor{secondblue}{RGB}{200,230,255}

\title{Extended Evaluation of SnowPole Detection for Machine-Perceivable Infrastructure for Nordic Winter Conditions: A Comparative Study of Object Detection Models}

\titlerunning{SnowPole Detection: Comparative YOLO Evaluation}
 \date{\today}
\author{
Durga Prasad Bavirisetti\inst{1,2} \and
Muhammad Ibne Rafiq\inst{3} \and
Shaira Tabassum\inst{1} \and
Gabriel Hanssen Kiss\inst{1} \and
Frank Lindseth\inst{1}
}

\authorrunning{D.P. Bavirisetti et al.}

\institute{
Dept. of Computer Science, NTNU, Trondheim, Norway \and
Dept. of Computer Science, University of Gävle, Gävle, Sweden \and
Dept. of Mathematics and Computer Science, Eindhoven University of Technology, Netherlands \\
\email{durga.prasad.bavirisetti@hig.se} \\
}

\begin{document}



\maketitle
\begin{abstract}
-RGB+ A will work as shown with infrarer and they are same because both gives spatial and geometric stuff
-see if dual branch is actually working and standard and then see how the current one is performing. USE NOT A SIMPLE LOSS BUT EXACTLY WHAT v9t would do. make it absolute same then run it, see online or smth
-repeat it for comb3 comb 4 comb 5 and attention or not then present it to DP SIR


-do repeat of v9t, v11n and comb3, 4 repeatedly on the same techniQUE OF 5 EXPERIMENTS.
-see if u can use merkle tree for speeding up MULTI-BRANCH to take into account for all the modalities?
-try other dataset for genelizability & ask the lab for it to provide other dataset
-but first the 5 main steps:
for doing experiments 2, 3, 5, and 7 (attention ablation, modality ablation, fusion architecture test, and Merkle tree/inference caching) with YOLOv9c as baseline, plus detailed instructions on how to do step 1 for YOLOv9c:

## Step 1: Baseline Setup on YOLOv9c

1. **Setup Environment and Clone YOLOv9 Repository**
   - Clone the YOLOv9 repo (e.g., https://github.com/WongKinYiu/yolov9 or maintained forks).
   - Install all dependencies from `requirements.txt`.
2. **Download Pretrained YOLOv9c Weights**
   - Download weights (e.g., yolov9-c.pt) for pretrained RGB model from official release.
3. **Prepare Dataset with RGBA inputs**
   - Organize your dataset folder with your RGB images and matched additional modality images (range/IR) for alpha.
   - Label annotations must remain consistent in COCO/YOLO format.
4. **Modify Input Layer to Accept RGBA (4 Channels)**
   - Adjust first convolutional input channel from 3 (RGB) to 4 or dual-branch for separate inputs.
   - Initialize RGB channels with pretrained weights, initialize A channel with zeros or small random weights.
5. **Train YOLOv9c with fused input**
   - Use training scripts, specifying your modified dataset and possibly longer epochs for convergence.
6. **Evaluate baseline performance with fused input (precision, recall, mAP, inference speed).**

***

## Step 2: Attention Module Ablation (Experiment 2)

- Add attention modules (e.g., Efficient Multi-scale Attention) after fusion layers or feature extraction in the backbone.
- Train the full model with attention and save metrics.
- Remove or bypass attention layers; retrain or fine-tune without them.
- Compare precision, recall, mAP, and inference time.
- This quantifies attention's contribution.

***

## Step 3: Modality Branch Ablation (Experiment 3)

- Train/test variants:
  - RGB only (original YOLO).
  - A channel only (range or infrared).
  - Combined RGB + A channels (fused or dual-branch).
- Keep everything else constant.
- Evaluate how much the additional modality boosts detection performance.

***

## Step 4: Fusion Architecture Ablation (Experiment 5)

- Swap fusion types to compare:
  - Simple concatenation of RGBA input.
  - Element-wise addition post feature extraction of separate branches.
  - Gated fusion where learnable weights modulate feature contributions.
- Ensure consistent training schedules.
- Measure trade-offs in precision, recall, mAP, and inference latency.

***

## Step 5: Merkle Tree or Tile-based Inference Caching (Experiment 7)

- Implement caching of unchanged spatial regions between frames/inputs using:
  - Tile-based difference checks.
  - Merkle or hash tree computation on input partitions.
- Skip inference for cached tiles and merge predictions for full frame.
- Benchmarks inference speed gains and any minor impacts on detection metrics.



All datasets, code, and trained models are publicly available to support reproducibility via our GitHub repository\footnote{\url{https://github.com/MuhammadIbneRafiq/Extended-evaluation-snowpole-lidar-dataset}} and the Mendeley dataset archive\footnote{\url{https://data.mendeley.com/datasets/tt6rbx7s3h/3}}.
\end{abstract}







\textbf{Keywords:} Snow Pole Detection, Object Detection, LiDAR Perception, YOLO Models, Multimodal Fusion, Localization, Nordic Winter Conditions

% \begin{multicols}{2} % Start two-column environment



\section{Introduction}

Robust detection and localization of roadside elements such as snow poles are vital for infrastructure safety and reliable autonomous vehicle operation during harsh Nordic winters. Snow poles serve as essential visual references when accumulated snow obscures lanes, curbs, and other markings. However, accurately identifying these objects is challenging due to unpredictable weather, occlusions, varying lighting conditions, and the inherent sparsity and noise of LiDAR data in real-world scenarios.

In prior work~\cite{bavirisetti2024pole}, we developed a YOLOv5s-based detection and geolocation pipeline on LiDAR signal images, establishing a baseline for this problem. Building on this, we released the SnowPole Detection dataset~\cite{bavirisetti2025snowpole}, which includes diverse annotated 360-degree LiDAR scans collected with an OS2-128 sensor mounted on an autonomous vehicle across mountainous roads, open terrain, and forested areas.

The raw point clouds were converted to panoramic LiDAR images across four modalities: Near-Infrared (Near-IR), Signal, Reflectivity, and Range. To support further experimentation, pseudo-color images were created by mapping these channels to RGB layers. Annotation was conducted in Roboflow and refined in CVAT to ensure high-quality labels. The dataset comprises 1,954 labeled samples split 70/20/10 for training, validation, and testing. Since modalities are spatially aligned, labels are consistent across channels, and the dataset is publicly available on Mendeley Data.

In this work, we present an extended version of the dataset~\cite{bavirisetti2025extended} (Mendeley Version 3), which includes the original modalities, the initial RGB combination, and five new pseudo-color configurations. While earlier results confirmed feasibility with single-channel inputs, key questions remained about how recent object detectors would perform across modalities and what trade-offs would arise in detection performance and latency.

Recent studies highlight that combining complementary sensor channels can improve detection robustness in adverse conditions~\cite{zafar2024realtime,nature2023fusion}. Motivated by this, we systematically evaluate YOLOv5s, YOLOv7-tiny, YOLOv8n, YOLOv9t, YOLOv10n, and YOLOv11n across all input configurations. Our analysis uses standard detection metrics (Precision, Recall, mAP@50, mAP@50--95) and measures inference latency on GPU and CPU hardware. To enable systematic comparison, we also introduce a composite \emph{Rank Score} metric that integrates detection accuracy and real-time performance. Our objective is to identify configurations offering the best trade-off between detection quality and inference efficiency for deployment in challenging Nordic winter environments.


\section{Related Work}

Recent years have seen major advances in object detection frameworks, sensor fusion, and benchmarking for real-time perception in autonomous driving and infrastructure monitoring. This section reviews work in LiDAR-based object detection, LiDAR imaging modalities for enhanced perception, and real-time inference optimization.

\subsection{LiDAR-Based Object Detection}

LiDAR sensors have become a core modality for autonomous perception, offering high-resolution spatial information robust to illumination changes. Early pipelines relied on voxelization and multi-view projections of point clouds, such as MV3D~\cite{chen2017multi} and VoxelNet~\cite{zhou2018voxelnet}. While accurate, their computational demands limited deployment on embedded systems.

To improve efficiency, PointPillars~\cite{lang2019pointpillars} and PV-RCNN~\cite{shi2020pv} introduced hybrid encoding schemes combining point- and grid-based features. These architectures achieved strong results on benchmarks like KITTI but were mainly designed for detecting large objects such as vehicles and pedestrians, rather than roadside infrastructure.

More recently, 2D CNN-based detection on LiDAR-derived range or pseudo-color images has gained traction. The YOLO series has been widely adopted for its balance of speed and accuracy. From YOLOv3~\cite{redmon2018yolov3} and YOLOv4~\cite{bochkovskiy2020yolov4} to YOLOv7’s trainable bag-of-freebies~\cite{wang2023yolov7} and YOLOv8’s anchor-free design~\cite{jocher2023ultralytics}, each iteration has introduced refinements targeting real-time inference. Benchmark studies have shown YOLO variants reaching over 60 mAP while sustaining more than 40 FPS on embedded platforms such as NVIDIA Jetson AGX~\cite{lazarevich2023yolobench}.

Transformer-based detectors, including Vision Transformers (ViTs)~\cite{dosovitskiy2020image} and Swin Transformers~\cite{liu2022swin}, have demonstrated state-of-the-art perception performance by modeling long-range dependencies, though their high computational demands currently limit adoption in latency-sensitive settings.


\subsection{LiDAR Imaging Modalities for Enhanced Perception}

Bavirisetti et al.~\cite{bavirisetti2024pole} introduced a LiDAR signal imaging-based pole detection framework that achieved promising detection accuracy in challenging conditions. However, their approach focused exclusively on a single-channel modality. In contrast, our present study offers a more comprehensive evaluation by analyzing detection performance across multiple LiDAR imaging modalities as well as pseudo-color representations created through intra-LiDAR fusion of different channels.

Recent studies have demonstrated that multi-spectral LiDAR data can significantly improve the object discrimination tasks. For example, Valentin et al.~\cite{vierhub2022study} showed that the effect of multispectral Lidar data on automated semantic segmentation in complex road environments. Building on these insights, our work systematically investigates how fusing Reflectance, Signal, and Near-Infrared modalities into pseudo-color images affects detection accuracy and real-time feasibility.

% ~\cite{bijelic2020seeing}
\subsection{Real-Time Inference and Benchmarking}

For practical deployment, detection models must balance accuracy with low latency. Two-stage detectors such as Faster R-CNN~\cite{ren2015faster} continue to be preferred in safety-critical scenarios due to their high precision. However, their runtime often limits suitability for embedded systems. In contrast, single-stage detectors such as YOLO algorithms are optimized for real-time inference, achieving throughput suitable for closed-loop path planning and obstacle avoidance~\cite{wang2023yolov7}.

Emerging frameworks such as EfficientDet~\cite{tan2020efficientdet} have proposed composite evaluation metrics that integrate accuracy and latency, enabling better comparison of model architectures. However, no standardized ranking method exists specifically for roadside object detection using LiDAR modalities. This motivates our work, which introduces a unified Rank Score metric that jointly considers detection accuracy and inference speed to guide model and modality selection for real-time applications.


\section{Methods and Experiments}

This section describes the dataset, experimental setup, training procedures, and inference protocols employed in this work. The objective was to conduct a rigorous and reproducible evaluation of recent YOLO architectures for snow pole detection across a diverse range of LiDAR imaging modalities and pseudo-color channel combinations. This section also presents the Composite Ranking Score formulation used to integrate detection performance and inference latency into a unified evaluation metric.

\subsection{Dataset Description}

In this work, we present an extended release of the SnowPole Detection dataset~\cite{bavirisetti2025extended} (Mendeley Dataset Version~3). Compared to the initial version, which included only the original single-channel modalities and one pseudo-color configuration, this updated dataset introduces five additional RGB combinations derived by mapping the single channels to different color representations. 

The dataset consists of 360-degree LiDAR images acquired using an OS2-128 high-resolution sensor mounted on an autonomous vehicle research platform~\cite{bavirisetti2025snowpole}. The data is split into training, validation, and test partitions according to a 70\%/20\%/10\% ratio to facilitate reproducibility and consistent evaluation. 

% \setcounter{figure}{0}  % reset so next figure is #1
\begin{figure}[htbp]

\renewcommand{\figurename}{Figure}
  \centering
  % Row 1
  \begin{subfigure}[b]{0.48\textwidth}
    \includegraphics[width=\textwidth]{ground_truth_reflec_1967.png}
    \caption{Reflectance}
    \label{fig:reflect_original}
  \end{subfigure}
  \hfill
  \begin{subfigure}[b]{0.48\textwidth}
    \includegraphics[width=\textwidth]{ground_truth_signal_1967.png}
    \caption{Signal}
    \label{fig:signal_original}
  \end{subfigure}
  % Row 2
  \begin{subfigure}[b]{0.48\textwidth}
    \includegraphics[width=\textwidth]{ground_truth_nearir_1967.png}
    \caption{Near-Infrared}
    \label{fig:nearir_original}
  \end{subfigure}
  \hfill
  \begin{subfigure}[b]{0.48\textwidth}
    
\includegraphics[width=\textwidth]{ground_truth_range_1967.png}
    \caption{Range}
    \label{fig:range_original}
  \end{subfigure}
  % Row 3
  \begin{subfigure}[b]{0.48\textwidth}
    \includegraphics[width=\textwidth]{ground_truth_combination1_1967.png}
    \caption{Combination 1}
    \label{fig:combination1_original}
  \end{subfigure}
  \hfill
  \begin{subfigure}[b]{0.48\textwidth}
    \includegraphics[width=\textwidth]{ground_truth_combined_color_1967.png}
    \caption{Combination 2}
    \label{fig:combination2_original}
  \end{subfigure}
  % Row 4
  \begin{subfigure}[b]{0.48\textwidth}
    \includegraphics[width=\textwidth]{ground_truth_combination3_1967.png}
    \caption{Combination 3}
    \label{fig:combination3_original}
  \end{subfigure}
  \hfill
  \begin{subfigure}[b]{0.48\textwidth}
    \includegraphics[width=\textwidth]{ground_truth_combination4_1967.png}
    \caption{Combination 4}
    \label{fig:combination4_original}
  \end{subfigure}
  % Row 5
  \begin{subfigure}[b]{0.48\textwidth}
    \includegraphics[width=\textwidth]{ground_truth_combination5_1967.png}
    \caption{Combination 5}
    \label{fig:combination5_original}
  \end{subfigure}
  \hfill
  \begin{subfigure}[b]{0.48\textwidth}
    \includegraphics[width=\textwidth]{ground_truth_combination6_1967.png}
    \caption{Combination 6}
    \label{fig:combination6_original}
  \end{subfigure}

  \caption{Overview of all single modalities and pseudo-color combinations avaialble in the dataset with the groundtruth poles\cite{bavirisetti2025extended}.}
  \label{figure:all_permutations_overview}
\end{figure}
As illustrated in Figure~\ref{figure:all_permutations_overview}, the dataset includes four single-modality inputs and six pseudo-color configurations created by combining these channels into RGB images:

\begin{itemize}
    \item \textbf{Reflectance}: Encodes surface reflectivity, which is helpful for distinguishing between materials.
    \item \textbf{Signal}: Captures intensity measurements from LiDAR returns.
    \item \textbf{Near-Infrared (Near-IR)}: Records near-infrared reflectance characteristics.
    \item \textbf{Range}: Provides distance information to estimate the location of detected poles.
    \item \textbf{Combination 1}: A pseudo-color image mapping Reflectance (R), Near-IR (G), and Signal (B) channels to the red, green, and blue channels, respectively.
    \item \textbf{Combination 2}: A pseudo-color image mapping  Reflectance(R), Signal (G), and Near-IR (B) channels to RGB.
    \item \textbf{Combination 3}: A pseudo-color image mapping Near-IR (R), Reflectance (G), and Signal (B).
    \item \textbf{Combination 4}: A pseudo-color image mapping Near-IR (R), Signal (G), and Reflectance (B).
    \item \textbf{Combination 5}: A pseudo-color image mapping Signal (R), Reflectance (G), and Near-IR (B).
    \item \textbf{Combination 6}: A pseudo-color image mapping Signal (R), Near-IR (G), and Reflectance (B).
\end{itemize}

This richer set of modalities and channel combinations enables a systematic investigation of how different input configurations impact detection performance and inference efficiency.



\subsection{Experimental Setup, Model Selection, Training, Validation, and Inference}

To ensure a rigorous and fair comparison, we evaluated six state-of-the-art YOLO architectures: YOLOv5s, YOLOv7-tiny, YOLOv8n, YOLOv9t, YOLOv10n, and YOLOv11n~\cite{ultralytics2025yolov5,wang2023yolov7,wang2024yolov9,jegham2023comprehensive}. These models were implemented using the Ultralytics framework and corresponding open-source repositories~\cite{ultralytics2025docs,ultralytics2024benchmark,ultralytics2024blog}. For consistency, we selected the smallest officially released variant of each YOLO version, ensuring all were optimized for real-time performance. This design isolates the impact of architectural innovations and incremental improvements across YOLO generations, allowing any observed performance differences to be attributed to the models themselves rather than variations in training configuration~\cite{ultralytics2024yolo11,jegham2025yolo}.

All models were trained using consistent hyperparameters, including identical input resolution, batch size, learning rate schedules, and data augmentation strategies. Training was conducted for 150 to 300 epochs depending on convergence behavior, with early stopping criteria enabled to mitigate overfitting and reduce unnecessary computational overhead~\cite{goodfellow2016deep}.Each YOLO variant was trained using the Ultralytics command-line interface (CLI) and the official configuration files, adhering to reproducible workflows and established best practices in object detection research~\cite{redmon2016you}. All scripts, configuration files, and training commands have been made publicly available to facilitate transparency and replicability.\footnote{\url{https://github.com/MuhammadIbneRafiq/Extended-evaluation-snowpole-lidar-dataset}}

The computational environment was standardized across all experiments and comprised a DELL workstation equipped with an Intel Core Ultra 7 (i7-13700HX) processor featuring 16 physical cores (8 performance and 8 efficiency cores) and 24 threads. The GPU used for training and inference was an NVIDIA RTX 2000 Ada Generation Laptop GPU with 8~GB GDDR6 VRAM and 3072 CUDA cores, leveraging the Ada Lovelace architecture for accelerated deep learning workloads. The system included 32~GB of DDR5 memory and a 1~TB NVMe PCIe Gen4 SSD to ensure efficient data loading and storage operations. The software environment was based on Python~3.10 and PyTorch~2.1 configured with CUDA~12.4, running on Windows~11 Pro for a stable and reproducible setup~\cite{jegham2023comprehensive,ultralytics2024blog}.

Validation employed the cross-validation procedures integrated into the Ultralytics YOLO framework, with default thresholds and Intersection-over-Union (IoU) parameters applied consistently across all experiments~\cite{bishop2006pattern}. This ensured an unbiased assessment of each model’s ability to generalize to unseen data partitions. For the held-out test set (197 images), key metrics—including mAP, Precision, Recall, and inference latency—were recorded to evaluate end-to-end detection performance comprehensively. 

Inference was conducted on both GPU and CPU hardware to evaluate real-time feasibility. Performance metrics were averaged over 10 runs per model. Additional qualitative inference was performed on single images using custom scripts, producing the visual examples presented in Figure~2. CPU inference was emphasized to assess suitability for edge deployment scenarios where GPU resources may be limited~\cite{lane2016deepx}. Furthermore, to assess the impact of input modality, all models were trained and evaluated across each individual modality (Reflectance, Signal, Near-Infrared) and six pseudo-color combinations. This systematic exploration allowed us to quantify how different YOLO architectures and input configurations influence detection accuracy and computational efficiency in practical deployment contexts~\cite{zhang2016color}.



\subsection{Composite Ranking Score for Model Evaluation}
\label{Ranking_Score}

To systematically evaluate and compare the performance of various YOLO architectures across different LiDAR-derived image modalities, we define a composite \textit{Rank Score} that aggregates multiple detection metrics into a single scalar value. This formulation enables the selection of model--modality pairs that optimally balance detection accuracy and inference performance, following practices inspired by prior studies combining accuracy and latency in benchmarking pipelines~\cite{bochkovskiy2020yolov4,lin2014microsoft,tan2020efficientdet}. In applications such as autonomous driving and robotics, object detection must balance predictive accuracy and computational latency. While metrics like mAP@50 and mAP@50--95 quantify detection quality, they do not fully capture real-time constraints. To address this, we define a rank-based scoring mechanism integrating detection performance and inference latency, providing a comprehensive basis for selecting models for edge deployment and safety-critical scenarios.



\subsubsection{Baseline Rank Score}

The baseline Rank Score combines four established detection metrics: mAP@50--95, mAP@50, Precision, and Recall. The weighting scheme prioritizes comprehensive detection accuracy while retaining direct measures of classification correctness and sensitivity:

\small{
\begin{equation}
\text{Rank Score}_{\text{baseline}} ={}&
  3\,\text{mAP}_{50\text{--}95}
  + 2\,\text{mAP}_{50}\\
&\quad+ \text{P}
  + \text{R}.
\end{equation}
}




The weights \((3, 2, 1, 1)\) were empirically chosen to reflect the relative importance of each metric in capturing overall detection quality. Specifically, \(\text{mAP}_{50\text{--}95}\) was assigned the highest weight (3) because it offers the most rigorous assessment of localization accuracy across diverse IoU thresholds, providing a comprehensive measure of model precision. \(\text{mAP}_{50}\) received a slightly lower weight (2), as it reflects performance under a more lenient yet widely used benchmark. Precision and Recall were given equal weights (1) to balance the model’s ability to minimize false positives (Precision) and maximize detection completeness (Recall). This weighting strategy ensures that models are evaluated not only for high average precision but also for reliable detection coverage and consistent classification across different scenarios.



\subsubsection{GPU-Based Rank score}

To incorporate inference efficiency, the Rank Score extends the baseline formulation with a latency penalty term proportional to the average GPU inference time in milliseconds:

\small{
\begin{equation}
\text{Rank Score}_{\text{GPU}} ={}&
  3\,\text{mAP}_{50\text{--}95}
  + 2\,\text{mAP}_{50}\\
&\quad+ \text{P}
  + \text{R}
  - \alpha\,\text{GPU}_{\text{ms}}.
\end{equation}
}


where $\alpha$ is a tunable hyperparameter controlling the impact of latency. In this study, we set $\alpha = 0.01$, penalizing every additional 10\,ms of GPU inference time by 0.1 points. This value was chosen to balance the trade-off between detection accuracy and computational efficiency such that latency has a measurable but not overwhelming influence on the composite score. The factor ensures that small variations in inference time do not disproportionately penalize high-accuracy models, while still promoting configurations suitable for real-time deployment scenarios.


\subsubsection{CPU-Based Rank score}

While GPU latency is critical for embedded and edge deployments, some applications—such as low-cost roadside units or traditional servers—rely primarily on CPU inference performance. To capture this scenario, we define an additional Rank Score that penalizes average CPU inference time (in milliseconds):

\small{
\begin{equation}
\text{Rank Score}_{\text{CPU}} ={}&
  3\,\text{mAP}_{50\text{--}95}
  + 2\,\text{mAP}_{50}\\
&\quad+ \text{P}
  + \text{R}
  - \alpha\,\text{CPU}_{\text{ms}}.
\end{equation}
}


Here, the coefficient $\alpha$ retains the same value ($\alpha=0.01$), implying that every additional 10\,ms of CPU inference time reduces the score by 0.1 points. This formulation supports ranking models in scenarios where GPU acceleration is unavailable or where CPU latency is a dominant deployment constraint.

In practice, the CPU-based Rank Score provides complementary insights alongside the GPU-adjusted score, helping practitioners select detection architectures that balance accuracy and computational feasibility under different hardware limitations.

This composite formulation offers several desirable properties: models with higher detection accuracy naturally receive higher scores, while models exhibiting higher inference latency incur a proportional penalty, promoting configurations suitable for real-time deployment. In addition, the formulation is designed to support practical scenarios where latency is a critical operational constraint, such as perception pipelines for autonomous vehicles or embedded edge AI platforms, ensuring that model selection balances both predictive performance and computational feasibility.




\section{Results and Analysis}

This section presents a comprehensive evaluation of the YOLO variants for snow pole detection across different LiDAR modalities and their combinations.

% \setcounter{figure}{0}  % reset so next figure is #1
\begin{figure}[htbp]
\centering

\begin{subfigure}[b]{0.9\textwidth}
    \includegraphics[width=\textwidth]{reflec.jpg}
    \caption{Reflectance}
    \label{fig:reflec}
\end{subfigure}

\begin{subfigure}[b]{0.9\textwidth}
    \includegraphics[width=\textwidth]{signal.jpg}
    \caption{Signal}
    \label{fig:signal}
\end{subfigure}

\begin{subfigure}[b]{0.9\textwidth}
    \includegraphics[width=\textwidth]{nearir.jpg}
    \caption{NearIR}
    \label{fig:nearir}
\end{subfigure}

\begin{subfigure}[b]{0.9\textwidth}
    \includegraphics[width=\textwidth]{Combination1.jpg}
    \caption{Combination 1 (R = Reflectance, G = NearIR, B = Signal)}
    \label{fig:comb1}
\end{subfigure}

\begin{subfigure}[b]{0.9\textwidth}
    \includegraphics[width=\textwidth]{Combination2.jpg}
    \caption{Combination 2 (R = Reflectance, G = Signal, B = NearIR)}
    \label{fig:comb2}
\end{subfigure}

\begin{subfigure}[b]{0.9\textwidth}
    \includegraphics[width=\textwidth]{Combination3.jpg}
    \caption{Combination 3 (R = NearIR, G = Reflectance, B = Signal)}
    \label{fig:comb3}
\end{subfigure}

\begin{subfigure}[b]{0.9\textwidth}
    \includegraphics[width=\textwidth]{Combination4.jpg}
    \caption{Combination 4 (R = NearIR, G = Signal, B = Reflectance)}
    \label{fig:comb4}
\end{subfigure}

\begin{subfigure}[b]{0.9\textwidth}
    \includegraphics[width=\textwidth]{Combination5.jpg}
    \caption{Combination 5 (R = Signal, G = Reflectance, B = NearIR)}
    \label{fig:comb5}
\end{subfigure}

\begin{subfigure}[b]{0.9\textwidth}
    \includegraphics[width=\textwidth]{Combination6.jpg}
    \caption{Combination 6 (R = Signal, G = NearIR, B = Reflectance)}
    \label{fig:comb6}
\end{subfigure}

\caption{ Example YOLOv11 Detection Results on Different LiDAR Modalities and Pseudo-Color Combinations.}
\label{fig:yolo_2}
\end{figure}

\subsection{Visual Quality Analysis}

Figures~\ref{fig:yolo_2} illustrate qualitative detection results across all input modalities and pseudo-color combinations. The single-modality outputs (Figures~\ref{fig:reflec},~\ref{fig:signal}, and~\ref{fig:nearir}) show clear differences in contrast and pole visibility. Reflectance imagery (Figure~\ref{fig:reflec}) provides strong edge definition, helpful for identifying pole contours against snow. Signal images (Figure~\ref{fig:signal}) enhance distant pole visibility under challenging lighting but can introduce background clutter. Near-Infrared imagery (Figure~\ref{fig:nearir}) offers lower structural contrast, resulting in less distinct pole boundaries, consistent with the lower mAP observed in quantitative evaluations. While inferencing on these image, the confidence threshold of the bounding box was taken as 0.4.

The pseudo-color combinations (Figures~\ref{fig:comb1}--\ref{fig:comb6}) highlight the benefits of multi-channel fusion. Combination~1 and Combination~2 (Figures~\ref{fig:comb1} and~\ref{fig:comb2}) show moderate improvements by merging complementary spectral information, though some noise remains. Combination~3 (Figure~\ref{fig:comb3}) and Combination~4 (Figure~\ref{fig:comb4}) consistently produce clearer pole delineation and reduced background clutter, enhancing detection across varying distances. Combination~5 (Figure~\ref{fig:comb5}) and Combination~6 (Figure~\ref{fig:comb6}) further improve contrast and separation but can suppress finer details compared to Combinations~3 and~4.

These visual trends align with the higher detection scores achieved by YOLOv9t and YOLOv11n in multi-channel configurations. The pseudo-color mappings effectively integrate information from Near-Infrared, Signal, Reflectance, and Range channels, demonstrating their utility for robust pole detection in snow-covered environments.

Overall, this analysis confirms that modality selection and fusion strategies significantly affect detection clarity and model confidence. These observations are consistent with the quantitative Rank Score evaluations and support adopting multi-channel combinations for deployment in real-world scenarios.





\subsection{Quantitative Analysis}
The preceding quantitative analysis provides detailed numerical insights into how detection accuracy and latency vary across models and configurations. To further aid interpretation of these results and support clear comparisons, we applied consistent visual encoding in the summary tables.

\begin{table}[h!]
    \centering
    \tiny
    \caption{Benchmarking of YOLO Architectures Across LiDAR Modalities}
\begin{tabular}{@{}llcccccccccc@{}}
\toprule
\textbf{Modality} & \textbf{Model} 
& \textbf{Precision} & \textbf{Recall} 
& \textbf{mAP50} & \textbf{mAP50–95} 
& \textbf{GPU (ms)} & \textbf{CPU (ms)} 
& \multicolumn{3}{c}{\textbf{Rank Score}} \\
\cmidrule(lr){9-11}
& & & & & & & & \textbf{Baseline} & \textbf{GPU} & \textbf{CPU} \\
\midrule


 \multirow{}{}{Reflectance}
& YOLO v5s 
& 0.851 & 0.734 & 0.797 & 0.320 
& 6.0 & 18.4 
& 4.139 & 4.079 & 3.955 \\
& YOLO v7-tiny 
& 0.867 & 0.788 & 0.854 & 0.354 
& 5.2 & 18.1 
& 4.425 & 4.373 & 4.244 \\
& YOLO v8n 
& \textit{0.877} & 0.801 & 0.859 & 0.390 
& \textbf{2.1} & \textit{17.8} 
& 4.566 & \textit{4.545} & 4.388 \\
& YOLO v9t 
& 0.861 & \textit{0.805} & \textit{0.871} & \textit{0.404} 
& 8.0 & 21.3 
& \textit{4.620} & 4.540 & \textit{4.407} \\
& YOLO v10n 
& 0.856 & 0.747 & 0.838 & 0.384 
& 4.8 & 20.1 
& 4.431 & 4.383 & 4.230 \\
& YOLO v11n 
& \textbf{0.884} & \textbf{0.824} & \textbf{0.886} & \textbf{0.405} 
& \textit{3.7} & \textbf{16.9} 
& \textbf{4.695} & \textbf{4.658} & \textbf{4.526} \\
\midrule
\multirow{}{}{Signal}
& YOLO v5s     
& \textit{0.888} & \textit{0.830} & 0.888 & 0.408 
& \textbf{2.8} & 17.9 
& 4.718 & 4.690 & 4.539 \\
& YOLO v7-tiny 
& 0.875 & 0.796 & 0.864 & 0.377 
& 8.4 & 19.6 
& 4.530 & 4.446 & 4.334 \\
& YOLO v8n     
& \textit{0.888} & 0.799 & 0.875 & \textbf{0.439} 
& \textit{3.8} & \textit{17.8} 
& 4.754 & 4.716 & 4.576 \\
& YOLO v9t     
& \textbf{0.894} & 0.813 & \textbf{0.901} & \textit{0.434} 
& 5.9 & 18.2 
& \textbf{4.811} & \textbf{4.752} & \textbf{4.629} \\
& YOLO v10n    
& 0.852 & 0.770 & 0.867 & 0.423 
& 5.0 & 19.7 
& 4.625 & 4.575 & 4.428 \\
& YOLO v11n    
& 0.850 & \textbf{0.843} & \textit{0.896} & 0.436 
& 4.4 & \textbf{17.6} 
& \textit{4.793} & \textit{4.749} & \textit{4.617} \\
\midrule


\multirow{}{}{Near-IR}
& YOLO v5s 
& 0.783 & 0.551 & 0.601 & 0.252 
& 4.9 & 18.7 
& 3.292 & 3.243 & 3.105 \\
& YOLO v7-tiny 
& 0.722 & 0.503 & 0.544 & 0.214 
& 5.1 & 19.2 
& 2.955 & 2.904 & 2.763 \\
& YOLO v8n 
& \textit{0.830} & 0.545 & 0.635 & 0.253 
& \textbf{3.7} & \textit{17.9} 
& 3.404 & \textit{3.367} & \textit{3.225} \\
& YOLO v9t 
& 0.782 & 0.552 & \textit{0.637} & \textit{0.268} 
& 5.6 & 19.8 
& \textit{3.412} & 3.356 & 3.214 \\
& YOLO v10n 
& 0.773 & \textit{0.554} & 0.617 & 0.253 
& \textit{4.8} & 19.5 
& 3.320 & 3.272 & 3.125 \\
& YOLO v11n 
& \textbf{0.832} & \textbf{0.582} & \textbf{0.669} & \textbf{0.275} 
& 6.2 & \textbf{17.1} 
& \textbf{3.577} & \textbf{3.515} & \textbf{3.406} \\
\midrule

% \midrule

\multirow{}{}{Combination 1}
& YOLO v5s 
& 0.889 & \textit{0.865} & \textbf{0.905} & 0.440 
& \textit{2.2} & 18.8 
& \textit{4.884} & \textbf{4.862} & \textit{4.696} \\
& YOLO v7-tiny 
& \textbf{0.905} & 0.864 & \textit{0.903} & 0.413 
& \textbf{2.1} & \textbf{17.4} 
& 4.814 & 4.793 & 4.640 \\
& YOLO v8n 
& 0.840 & 0.848 & \textbf{0.905} & 0.442 
& 8.4 & 18.6 
& 4.824 & 4.740 & 4.638 \\
& YOLO v9t 
& 0.840 & \textbf{0.873} & 0.902 & \textit{0.451} 
& 5.0 & 19.2 
& 4.870 & 4.820 & 4.678 \\
& YOLO v10n 
& \textit{0.899} & 0.768 & 0.883 & 0.445 
& 6.3 & 18.7 
& 4.768 & 4.705 & 4.581 \\
& YOLO v11n 
& 0.855 & 0.856 & 0.899 & \textbf{0.461} 
& 4.0 & \textit{17.9} 
& \textbf{4.892} & \textit{4.852} & \textbf{4.713} \\

\midrule
\multirow{}{}{Combination 2}
& YOLO v5s 
& 0.892 & \textbf{0.883} & \textit{0.908} & 0.440 
& 5.7 & 19.7 
& \textit{4.911} & \textit{4.854} & \textit{4.714} \\
& YOLO v7-tiny 
& \textbf{0.929} & \textit{0.868} & 0.907 & 0.438 
& 7.6 & 18.7 
& \textbf{4.925} & 4.849 & \textbf{4.738} \\
& YOLO v8n 
& 0.889 & 0.829 & \textbf{0.915} & \textbf{0.451} 
& \textbf{3.7} & \textbf{17.2} 
& 4.901 & \textbf{4.864} & 4.729 \\
& YOLO v9t 
& 0.875 & 0.838 & 0.889 & 0.441 
& 11.2 & 19.4 
& 4.814 & 4.702 & 4.620 \\
& YOLO v10n 
& \textit{0.893} & 0.787 & 0.891 & 0.442 
& \textit{5.2} & 18.3 
& 4.788 & 4.736 & 4.605 \\
& YOLO v11n 
& 0.881 & 0.843 & 0.905 & \textit{0.449} 
& \textbf{3.7} & \textit{17.4} 
& \textit{4.881} & 4.844 & 4.707 \\


\midrule
\multirow{}{}{Combination 3}
& YOLO v5s 
& 0.892 & \textbf{0.866} & \textit{0.910} & 0.434 
& \textit{2.2} & \textit{17.6}
& 4.880 & 4.858 & 4.710 \\
& YOLO v7-tiny 
& \textbf{0.905} & 0.840 & 0.890 & 0.410 
& \textbf{2.0} & \textbf{17.5} 
& 4.755 & 4.735 & 4.580 \\
& YOLO v8n 
& \textit{0.896} & 0.820 & 0.905 & 0.452 
& 3.1 & 18.0 
& 4.882 & 4.851 & 4.702 \\
& YOLO v9t 
& 0.895 & \textit{0.863} & \textbf{0.921} & \textit{0.454} 
& 4.3 & 19.0 
& \textbf{4.962} & \textbf{4.919} & \cellcolor{secondblue}\textbf{4.772} \\
& YOLO v10n 
& 0.889 & 0.810 & 0.903 & 0.437 
& 4.6 & 18.5 
& 4.816 & 4.770 & 4.631 \\
& YOLO v11n 
& 0.852 & \textbf{0.866} & 0.904 & \textbf{0.462} 
& 2.9 & 17.8 
& \textit{4.912} & \textit{4.883} & \textit{4.734} \\



\midrule
\multirow{}{}{Combination 4}
& YOLO v5s 
& \textbf{0.903} & \textit{0.864} & 0.906 & 0.438 
& \textit{3.6} & 17.9 
& 4.893 & \textit{4.857} & 4.714 \\
& YOLO v7-tiny 
& \textit{0.899} & 0.858 & 0.901 & 0.410 
& \textbf{2.1} & \textit{17.6} 
& 4.789 & 4.768 & 4.613 \\
& YOLO v8n 
& 0.859 & 0.865 & 0.904 & 0.448 
& 3.9 & 17.9 
& 4.876 & 4.837 & 4.697 \\
& YOLO v9t 
& 0.898 & \textbf{0.866} & \textbf{0.916} & \textbf{0.463} 
& 5.9 & 19.0 
& \cellcolor{bestgreen}\textbf{4.985} & \cellcolor{secondblue}\textbf{4.926} & \cellcolor{bestgreen}\textbf{4.795} \\
& YOLO v10n 
& 0.895 & 0.777 & 0.874 & 0.435 
& \textit{3.7} & 18.0 
& 4.725 & 4.688 & 4.545 \\
& YOLO v11n 
& 0.887 & \textbf{0.866} & \textit{0.914} & \textit{0.461} 
& 3.8 & \textbf{17.2} 
& \textit{4.946} & \cellcolor{secondblue}\textbf{4.926} & \textit{4.792} \\



  \midrule
\multirow{}{}{Combination 5}
& YOLO v5s 
& \textit{0.915} & \textit{0.863} & \textit{0.912} & 0.436 
& \textbf{2.1} & 17.9 
& 4.910 & 4.889 & 4.731 \\
& YOLO v7-tiny 
& 0.862 & \textbf{0.875} & 0.891 & 0.420 
& \textit{2.2} & \textit{17.5} 
& 4.779 & 4.757 & 4.604 \\
& YOLO v8n 
& \textbf{0.950} & 0.811 & \textbf{0.927} & 0.450 
& 3.1 & 17.8 
& \cellcolor{secondblue}\textbf{4.965} & \cellcolor{bestgreen}\textbf{4.934} & \textbf{4.787} \\
& YOLO v9t 
& 0.885 & 0.858 & 0.910 & \textbf{0.459} 
& 4.0 & 18.9 
& \textit{4.940} & \textit{4.900} & \textit{4.751} \\
& YOLO v10n 
& \textit{0.923} & 0.749 & 0.879 & 0.449 
& 3.5 & 18.2 
& 4.777 & 4.742 & 4.595 \\
& YOLO v11n 
& 0.878 & 0.857 & 0.899 & \textit{0.453} 
& 3.0 & \textbf{17.3} 
& 4.892 & \textit{4.862} & 4.719 \\

\midrule
\multirow{}{}{Combination 6}
& YOLO v5s 
& \textit{0.917} & 0.854 & 0.911 & 0.429 
& 3.5 & 17.6 
& \textit{4.880} & \textit{4.845} & \textit{4.704} \\
& YOLO v7-tiny 
& \textbf{0.918} & \textbf{0.867} & \textit{0.912} & 0.417 
& \textbf{2.2} & \textit{17.4} 
& 4.860 & 4.838 & 4.686 \\
& YOLO v8n 
& 0.876 & 0.840 & 0.902 & 0.449 
& \textit{3.1} & 17.7 
& 4.867 & 4.836 & \textit{4.690} \\
& YOLO v9t 
& 0.872 & \textit{0.862} & \textbf{0.917} & \textbf{0.461} 
& 5.4 & 18.8 
& \textbf{4.951} & \textbf{4.897} & \textbf{4.763} \\
& YOLO v10n 
& \textit{0.901} & 0.762 & 0.884 & 0.443 
& 4.8 & 18.1 
& 4.760 & 4.712 & 4.579 \\
& YOLO v11n 
& 0.858 & 0.823 & 0.889 & \textit{0.453} 
& 6.6 & \textbf{17.1} 
& 4.818 & \textit{4.752} & 4.647 \\

\bottomrule   
    \end{tabular}   
    \label{tab:yolo_results}
\end{table}

\subsubsection{Visual Encoding of Ranking Results}

To improve clarity and facilitate direct comparison across configurations, Table~\ref{tab:yolo_results} employs multiple visual annotations that highlight high-performing models. The conventions are summarized as follows:

\begin{itemize}
    \item \cellcolor{bestgreen}\textbf{Green shading} indicates the \emph{globally highest Rank Score} observed across all modalities and combinations. This helps quickly identify which model-modality pair achieved the most favorable trade-off between detection accuracy and latency overall.
    \item \cellcolor{secondblue}\textbf{Blue shading} denotes the \emph{second-highest global Rank Score} in the entire benchmark. This enables practitioners to see strong alternative candidates that are near the top in aggregate performance.
    \item \textbf{Bold font} highlights the \emph{best Rank Score within each individual modality or pseudo-color combination}. This local best helps identify which model is optimal when restricting deployment to a specific sensor configuration.
    \item \textit{Italic font} marks the \emph{second-best Rank Score within the same modality or combination}. This convention supports nuanced interpretation, showing strong alternatives when, for example, slight variations in mAP or latency may make the second-best preferable in a given scenario.
\end{itemize}

By combining these visual encodings—color shading, boldface, and italics—the table offers an intuitive and multi-layered comparison of detection accuracy and computational efficiency. This design facilitates the identification of both top-ranked configurations and strong secondary options for various deployment constraints.

\subsubsection{Detection Accuracy and mAP Analysis:}


Across individual modalities and pseudo-color combinations, \textbf{YOLOv9t}, \textbf{YOLOv8n}, and \textbf{YOLOv11n} consistently achieved the highest mAP$_{50}$ and mAP$_{50\text{--}95}$ values, demonstrating strong detection accuracy under diverse input conditions:

\begin{itemize}
    \item \textbf{YOLOv9t} achieved the highest overall mAP in many configurations, with mAP$_{50}$ reaching up to 0.921 (Combination~3) and mAP$_{50--95}$ as high as 0.463 (Combination~4). These high scores indicate excellent localization precision and robustness across a range of IoU thresholds, which is particularly important for reliably identifying snow poles under challenging conditions. YOLOv9t frequently ranked first among all models in both single-channel and pseudo-color modalities.
    
    \item \textbf{YOLOv8n} consistently delivered strong mAP performance across combinations, achieving mAP$_{50}$ up to 0.927 (Combination~5) and mAP$_{50--95}$ frequently above 0.44, demonstrating a solid balance of detection accuracy and real-time feasibility. Notably, in Combinations~2 and~5, YOLOv8n achieved results comparable to or slightly exceeding YOLOv9t in mAP$_{50}$, highlighting its competitive precision.
    
    \item \textbf{YOLOv11n} maintained high detection accuracy across most scenarios, with mAP$_{50}$ reaching 0.914--0.917 in several configurations (Combinations~4 and~6) and mAP$_{50--95}$ up to 0.462. Although its mAP$_{50}$ was slightly below YOLOv9t on average, YOLOv11n consistently ranked among the top-performing models, confirming that its architectural improvements contribute to robust detection.
    
    \item \textbf{YOLOv7-tiny} and \textbf{YOLOv5s} showed moderate mAP performance relative to the larger models. While their mAP$_{50}$ values were competitive—often exceeding 0.88--0.90 in the best configurations—they generally lagged behind YOLOv9t and YOLOv8n in mAP$_{50--95}$, with values between 0.410--0.440. This suggests that, although they can detect most poles reliably at standard IoU thresholds, their localization accuracy across stricter IoU ranges is more limited.
\end{itemize}

Overall, YOLOv9t emerged as the most accurate detector in terms of both mAP$_{50}$ and mAP$_{50--95}$, confirming its suitability for applications where the highest detection precision is required. YOLOv8n and YOLOv11n provided consistently strong mAP results across all modalities, offering compelling alternatives with excellent trade-offs between accuracy and efficiency.



\subsubsection{Precision and Recall Analysis}

Precision values across all evaluated models and input configurations remained consistently high, ranging from approximately 0.72 to 0.95. \textbf{YOLOv8n} demonstrated the highest observed Precision of 0.950 in Combination~5, confirming its excellent capability to produce reliable detections with minimal false positives. Similarly, \textbf{YOLOv7-tiny} and \textbf{YOLOv5s} often achieved Precision exceeding 0.90 in several pseudo-color configurations, such as Combinations~2 and~6, showing that even lightweight models can deliver strong precision when paired with well-chosen modalities. This high Precision is valuable in operational scenarios where reducing false alarms is critical, for example in automated alerting systems for roadside infrastructure monitoring.

Recall performance also remained robust across all models and modalities. The majority of configurations achieved Recall above 0.80, indicating strong sensitivity to detecting true positives. Notably, \textbf{YOLOv11n} and \textbf{YOLOv9t} consistently achieved Recall values in the range of 0.84 to 0.87, demonstrating their effectiveness in capturing nearly all instances of snow poles even in more challenging multimodal inputs. \textbf{YOLOv7-tiny} reached the highest Recall observed in the experiments, with a value of 0.875 in Combination~5, underscoring its suitability for applications where maximizing detection coverage is prioritized.

Overall, the analysis confirms that all tested YOLO architectures are capable of maintaining a strong balance between high Precision and high Recall across single-channel and pseudo-color LiDAR modalities. This balanced performance ensures that detection pipelines can reliably identify snow poles without a significant trade-off between minimizing false positives and capturing true positives, which is essential for real-time autonomous driving and winter infrastructure applications.



\subsubsection{Latency Analysis}

Inference latency varied notably across architectures and hardware platforms, underscoring the importance of selecting models that balance detection performance with responsiveness:

\begin{itemize}
    \item \textbf{GPU Inference}: Smaller models such as \textbf{YOLOv5s} and \textbf{YOLOv7-tiny} consistently achieved the fastest GPU inference times, typically between 2.0\,ms and 3.0\,ms across all modalities and combinations. This extremely low latency is highly advantageous for real-time deployments, including embedded perception systems in autonomous vehicles or roadside monitoring units that require rapid feedback. In contrast, larger models like \textbf{YOLOv9t} and \textbf{YOLOv10n} showed GPU latency rising up to approximately 11.2\,ms in the most computationally demanding configurations. While these runtimes still fall within the feasible range for near-real-time operation, they demand careful consideration in systems with strict frame rate or processing budgets.
    
    \item \textbf{CPU Inference}: As expected, CPU inference times were significantly higher across all models, ranging approximately between 17.0\,ms and 19.8\,ms. Notably, \textbf{YOLOv11n} consistently achieved the lowest CPU latency across most configurations (down to 17.1\,ms), reflecting its design optimizations for resource-constrained environments. This efficiency makes YOLOv11n particularly suitable for embedded or edge deployments that rely primarily on CPU execution without GPU acceleration. Conversely, heavier models such as \textbf{YOLOv9t} and \textbf{YOLOv10n} showed slightly higher CPU latencies, although the differences remained within a narrow margin of a few milliseconds.
\end{itemize}

These latency profiles directly influenced the GPU and CPU Rank Scores reported in Table~\ref{tab:yolo_results}. While \textbf{YOLOv9t} and \textbf{YOLOv11n} provided high detection accuracy, the trade-off was modestly increased latency compared to lighter variants like \textbf{YOLOv5s} and \textbf{YOLOv7-tiny}. Practitioners should weigh these considerations carefully: even small increases in inference time can have cumulative impacts on overall system throughput, especially in time-critical or multi-threaded environments. Selecting a model therefore requires balancing detection performance against the latency constraints of the deployment scenario.


\subsubsection{Composite Ranking and Trade-offs}

To facilitate systematic comparison, we introduced a composite Rank Score metric that integrates mAP$_{50\text{--}95}$, mAP$_{50}$, Precision, Recall, and latency penalties. This formulation enables clear differentiation of models that provide the optimal compromise between detection accuracy and real-time suitability across diverse LiDAR modalities and pseudo-color combinations. For example:

\begin{itemize}
    \item \textbf{YOLOv9t} consistently achieved the highest Baseline Rank Scores across several configurations, particularly in Combinations~3 and~4, where mAP$_{50}$ exceeded 0.92 and mAP$_{50\text{--}95}$ reached 0.463. These outstanding detection metrics confirm YOLOv9t as the most accurate model overall. However, it incurred moderate GPU and CPU Rank Score penalties due to higher inference latency relative to lighter variants.
    
    \item \textbf{YOLOv11n} frequently delivered the best CPU Rank Scores and among the top GPU Rank Scores. Its combination of high Precision and Recall (e.g., Recall up to 0.866) with the fastest CPU inference times (as low as 17.1\,ms) makes it an excellent choice for resource-constrained or CPU-only deployments where consistent detection accuracy must be balanced with efficient execution.
    
    \item \textbf{YOLOv8n} demonstrated strong Precision (up to 0.950) and consistently high mAP across multiple configurations while maintaining low GPU latency. As a result, it achieved the highest GPU Rank Scores in several cases, such as Combination~5 and Combination~3, making it particularly well-suited for GPU-accelerated real-time applications that require minimal processing delays without compromising accuracy.
    
    \item \textbf{YOLOv5s} and \textbf{YOLOv7-tiny} achieved the fastest inference on GPU (down to 2.0--2.2\,ms) and still delivered respectable mAP and Recall values, which led to competitive GPU Rank Scores despite their relatively simpler architectures. These models are therefore excellent candidates when ultra-low latency is the dominant requirement.
\end{itemize}

This structured analysis enables practitioners to identify the most appropriate model for their specific deployment constraints. In summary, \textbf{YOLOv9t} provides unmatched detection accuracy across modalities and is recommended in scenarios where maximizing mAP performance takes priority over latency constraints. \textbf{YOLOv11n} demonstrates the most consistent balance between detection metrics and computational efficiency, making it particularly attractive for CPU-limited environments requiring reliable real-time operation. Finally, \textbf{YOLOv8n}, \textbf{YOLOv7-tiny}, and \textbf{YOLOv5s} represent highly capable alternatives for GPU-based deployments where rapid inference is critical for maintaining system responsiveness in Nordic winter conditions.

\subsubsection{Recommended Modalities for Snow Pole Detection}

Based on the comprehensive evaluation of detection accuracy, precision, recall, and inference latency, several key observations emerge regarding the optimal choice of input modalities and pseudo-color combinations for snow pole detection in Nordic winter conditions:

\begin{itemize}
    \item \textbf{Pseudo-color combinations consistently outperformed single modalities.} Configurations combining multiple LiDAR channels---particularly Combinations~3, 4, and 5---produced the highest Baseline, GPU, and CPU Rank Scores across nearly all YOLO architectures. These combinations leverage complementary information from reflectance, signal strength, and near-infrared intensity to improve feature discriminability, resulting in higher mAP and more robust detections.
    
    \item \textbf{Combination~4 and Combination~5} were the most effective overall. These configurations yielded the top Baseline Rank Scores (often exceeding 4.9) and consistently high mAP$_{50}$ values (up to 0.927) with minimal latency penalties. This demonstrates that fusing reflectance and signal information with near-infrared intensities enhances model confidence and reduces both false negatives and false positives.
    
    \item \textbf{Single-modality inputs remain viable but show performance limitations.} While Reflectance and Signal alone still achieved strong detection metrics, they lagged behind combinations in Rank Scores, especially for YOLOv9t and YOLOv11n. Near-IR alone was the weakest modality in terms of precision, recall, and mAP.

\end{itemize}

Overall, multimodal LiDAR representations should be prioritized to ensure high-precision, reliable snow pole detection under harsh Nordic winter conditions.


\section{Conclusion}

This study presented an extensive evaluation of recent YOLO object detection architectures applied to the SnowPole Detection dataset, covering diverse LiDAR imaging modalities and pseudo-color combinations. By introducing a composite Rank Score metric that integrates detection accuracy and inference latency, we provided a principled framework for comparing models in scenarios requiring both high precision and real-time feasibility.

Our results show that YOLOv9t achieves the highest detection accuracy across most configurations, while YOLOv11n offers the best trade-off between accuracy and computational efficiency, particularly for CPU-based deployment. YOLOv7-tiny and YOLOv5s remain compelling choices when extremely low GPU latency is required.

The findings of this work contribute practical guidance for selecting object detection pipelines tailored to harsh Nordic winter conditions, supporting robust infrastructure monitoring and autonomous vehicle perception. Future research will explore further optimization of model architectures and the integration of temporal fusion strategies to enhance detection stability under challenging environmental variability.





\section*{Acknowledgments}
This research was supported by the Norwegian Research Council under the project "Machine Sensible Infrastructure under Nordic Conditions" (Project Number 333875). We also thank the NTNU NAPLab team for their support during data collection and vehicle setup.





\bibliography{references}
\bibliographystyle{plain}  % or your preferred style% 

\end{document}